\documentclass[12pt]{article}
\usepackage[utf8]{inputenc}
\usepackage[T1]{fontenc}
\usepackage{hyperref}
\title{Makolkol Language Resources for K-12 Curriculum Development}
\author{Compiled by Claude Sonnet 5.5 (claude-sonnet-5-5) for Andr\'as Kornai}
\date{}
\begin{document}
\maketitle

\section{Introduction}
Makolkol (ISO 639-3 zmh; Glottocode mako1250) was spoken on the Gazelle Peninsula, East New Britain, Papua New Guinea. It is very poorly documented and probably no longer spoken. Reported figures are 7 speakers in 1988, with use confined to the village of Gunapeo and a shift toward Tok Pisin and Meramera. Its classification is uncertain, and sources differ on whether it is related to the Baining languages.

\section{Primary Educational Resources}
None found. No school materials exist in the sources found in the searches.

\section{Linguistic and Reference Materials}
\begin{itemize}
\item Wikipedia, ``Makolkol language''. It mentions a 174-word list (Rosensteel 1988), the main documentation. \url{https://en.wikipedia.org/wiki/Makolkol_language}
\item Wikipedia, ``Baining languages'', for context on the neighbouring family. \url{https://en.wikipedia.org/wiki/Baining_languages}
\item Wikipedia, ``East New Britain languages''. \url{https://en.wikipedia.org/wiki/East_New_Britain_languages}
\end{itemize}

\section{Religious and Community Resources}
None found for Makolkol. For the neighbouring Baining people: \url{https://en.wikipedia.org/wiki/Baining_people}

\section{Digital and Audio Resources}
\begin{itemize}
\item Endangered Languages Project entry. \url{https://endangeredlanguages.com/lang/2268}
\item Verbix endangered-languages page. \url{https://docs.verbix.com/EndangeredLanguages/Makolkol}
\end{itemize}

\section{Curriculum Development Context}
The documentation is limited to a short wordlist, and no community of speakers is known. A curriculum in Makolkol is not feasible with the material found. At most, a cultural-heritage unit could be written in the local language of schooling (Tok Pisin or English), and any such work would need to involve local communities.

\end{document}